\documentclass[11pt]{article}

\usepackage[a4paper,margin=29mm]{geometry}
\usepackage[T1]{fontenc}
\usepackage[utf8]{inputenc}
\usepackage{lmodern,microtype}
\usepackage{amsmath,amssymb,amsthm,mathtools}
\usepackage{booktabs}
\usepackage[hidelinks]{hyperref}
\usepackage[nameinlink,noabbrev]{cleveref}

\newtheorem{theorem}{Theorem}[section]
\newtheorem{lemma}[theorem]{Lemma}
\newtheorem{proposition}[theorem]{Proposition}
\newtheorem{corollary}[theorem]{Corollary}
\theoremstyle{remark}
\newtheorem{remark}[theorem]{Remark}

\newcommand{\Z}{\mathbb Z}
\newcommand{\F}{\mathbb F}
\newcommand{\N}{\mathbb N}
\newcommand{\R}{\mathbb R}
\newcommand{\leg}[2]{\left(\frac{#1}{#2}\right)}
\newcommand{\pp}{\operatorname{pq}}
\newcommand{\ttq}{\operatorname{tq}}
\newcommand{\B}{\mathcal B}
\newcommand{\norm}[1]{\left\lVert #1\right\rVert}

\hypersetup{
  pdftitle={Integer quadruples and primitive three-square representations},
  pdfauthor={Nicolo Pandolfo},
  pdfsubject={Conventional proofs and Lean-certified arithmetic for OEIS A354766 and A278085},
  pdfkeywords={integer quadruples, primitive three squares, lattice neighbors, multiplicativity, natural density}
}

\title{\textbf{Integer quadruples and primitive three-square representations}\\
\large Product formulas, equality cases, and density zero for A354766/A278085}
\author{Nicol\`o Pandolfo}
\date{October 23, 2026\\\small Draft for independent human mathematical review}

\begin{document}
\maketitle

\noindent\textbf{Status and scope.}
This is a conventional mathematical proof, with its standard inputs stated
explicitly. A companion Lean development also verifies the representation
formula and the arithmetic, equality, and density conclusions without proof
admissions. Its final dependency audits use only the standard axioms
\texttt{propext}, \texttt{Classical.choice}, and \texttt{Quot.sound}.
Independent external human mathematical review is pending.

\begin{abstract}
We count ordered integer quadruples whose sum is a positive integer $n$ and
whose sum of squares is $n^2$. A Hadamard transformation gives a precise
primitive-counting bridge, including the different normalizations for odd
$n$ and for $n\equiv2\pmod4$. We prove the needed primitive three-square
product by an explicit lattice-neighbor double count. The only
non-elementary standard inputs are Minkowski's convex-body theorem and
quadratic reciprocity; no three-square counting formula is assumed. The
resulting normalized primitive and total counts are multiplicative, with
explicit prime-power values. We characterize exactly when the total count
equals $n$ and prove that this equality set has natural density zero by a
finite arithmetic sieve. The sieve requires no theorem about reciprocal
primes in an arithmetic progression.
\end{abstract}

\section{Definitions and proof dependencies}

For $n\geq1$, let $\mathcal Q_n$ be the set of ordered quadruples
$(h,i,j,k)\in\Z^4$ satisfying
\begin{equation}\label{eq:quad}
 h+i+j+k=n,\qquad h^2+i^2+j^2+k^2=n^2.
\end{equation}
Signs and coordinate orders are counted separately. Define the content of a
quadruple to be $\gcd(|h|,|i|,|j|,|k|)$; a quadruple is \emph{primitive}
when its content is one. Put
\[
 \ttq(n)=|\mathcal Q_n|,\qquad
 \pp(n)=|\{q\in\mathcal Q_n:q\text{ is primitive}\}|.
\]
These sets are finite: every coordinate has absolute value at most $n$.
Following the normalization of the OEIS pair~\cite{oeis-total,oeis-primitive},
put initially as rational numbers
\begin{equation}\label{eq:normalized}
 T(n)=\frac{\ttq(n)}4,\qquad P(n)=\frac{\pp(n)}4.
\end{equation}
We will prove divisibility of both raw counts by four before using these as
integer-valued arithmetic functions.

For a positive integer $M$, let $r_3^{\mathrm{prim}}(M)$ count ordered
integer triples $(x,y,z)$ with $x^2+y^2+z^2=M$ and
$\gcd(|x|,|y|,|z|)=1$. In particular, define
\[
 F(m)=r_3^{\mathrm{prim}}(3m^2)\qquad(m\geq1).
\]
All these sets are finite by the same coordinate bound. For an odd prime
$p$, let $\chi_p=\leg{-3}{p}$ be the Legendre symbol, with value zero when
$p=3$.

The lattice construction below is the classical neighbor construction of
Kneser~\cite{kneser}; Voight~\cite{voight}, Proposition~3.13 and
Corollary~3.15, provide background. We prove the particular constructions,
low-rank classification, and representation double count needed here.
Cooper--Hirschhorn~\cite{cooper-hirschhorn} give an independent published
cross-check of the product formula, but their representation theorem is not
an input to this proof.

The standard inputs used without reproof are Minkowski's theorem in the
following form: a centrally symmetric convex open set in $\R^r$ of volume
strictly greater than $2^r$ times a lattice's covolume contains a nonzero
lattice vector~\cite{minkowski}; and quadratic reciprocity with the
supplementary law for $-1$~\cite{gauss}. We also use elementary integer
factorization, the Chinese remainder theorem, finite-dimensional linear
algebra, and finite sums and products. No additional representation formula
or analytic assertion about primes is assumed.

\section{The corrected Hadamard counting bridge}

For a quadruple satisfying \eqref{eq:quad}, put
\[
 x=h+i-j-k,\quad y=h-i+j-k,\quad z=h-i-j+k.
\]
Expansion gives
\[
 x^2+y^2+z^2+(h+i+j+k)^2=4(h^2+i^2+j^2+k^2),
\]
so $x^2+y^2+z^2=3n^2$. Conversely, the inverse expressions are
\begin{equation}\label{eq:inverse}
 \begin{aligned}
 4h&=n+x+y+z,&4i&=n+x-y-z,\\
 4j&=n-x+y-z,&4k&=n-x-y+z.
 \end{aligned}
\end{equation}
Whenever their right sides are divisible by four, these expressions give
an integer quadruple satisfying \eqref{eq:quad}. Indeed, their sum is $4n$;
the same Hadamard norm identity gives the sum of squares $n^2$. Thus these
linear maps are inverse bijections after the integrality condition is imposed.

\begin{lemma}[A common divisor]\label[lemma]{lem:valuation}
If $d,m$ are positive integers and $d^2\mid3m^2$, then $d\mid m$.
\end{lemma}
\begin{proof}
For every prime $p$, write $a=v_p(d)$ and $b=v_p(m)$. If $p\neq3$,
the divisibility gives $2a\leq2b$. For $p=3$ it gives
$2a\leq1+2b$, which also implies $a\leq b$ because both are integers.
This proves the assertion prime by prime.
\end{proof}

The primary definition of A278085 uses primitive quadruples of sum $3n$
and squared norm $3n^2$. Its equivalence with $P(n)$ is explicit: the map
\[
 (h,i,j,k)\longmapsto(n-h,n-i,n-j,n-k)
\]
is an involution interchanging these quadruples with those in $\mathcal Q_n$.
Expansion proves the two sum and norm identities. It preserves content:
the content of a quadruple in $\mathcal Q_n$ divides $n$ by its sum, while
the content $d$ of a quadruple in the other family satisfies
$d^2\mid3n^2$, hence $d\mid n$ by \cref{lem:valuation}. In each direction
a common divisor therefore divides the complemented coordinates as well.
Thus the primitive subsets correspond bijectively.

\begin{theorem}[Primitive counting bridge]\label[theorem]{thm:bridge}
For every $n\geq1$,
\begin{equation}\label{eq:bridge}
 \pp(n)=
 \begin{cases}
 \frac12F(n),&n\text{ odd},\\[1mm]
 \frac12F(n/2),&n\equiv2\pmod4,\\[1mm]
 0,&4\mid n.
 \end{cases}
\end{equation}
\end{theorem}

\begin{proof}
Suppose first that $n$ is odd. A square is $0,1$, or $4$ modulo eight,
and three squares summing to $3n^2\equiv3\pmod8$ must all be one modulo
eight. Thus $x,y,z$ are odd. The four numerators in \eqref{eq:inverse} are
congruent modulo four, so integrality is equivalent to
\[
 n+x+y+z\equiv0\pmod4.
\]
Pair a triple with its negative. If $S=x+y+z$, then $n,S$ are odd, and
$n+S$ and $n-S$ are even residues whose sum is $2\pmod4$; exactly one is
zero modulo four. No triple is its own negative because its norm is positive.
Hence exactly one member of every sign pair gives an integer quadruple.

Primitivity is preserved for each integral inverse. A common divisor of the
quadruple divides $x,y,z$. Conversely, the gcd $d$ of $x,y,z$ is odd and
satisfies $d^2\mid3n^2$, so $d\mid n$ by \cref{lem:valuation}.
Since four is invertible modulo $d$, \eqref{eq:inverse} implies that $d$
divides all four coordinates. This proves the first line of \eqref{eq:bridge}.

Now let $n=2m$ with $m$ odd. From $x^2+y^2+z^2=12m^2\equiv0\pmod4$,
all three coordinates are even: the number of odd coordinates would otherwise
be one, two, or three, none of which gives zero modulo four. Write
$x=2X,y=2Y,z=2Z$. Then $X^2+Y^2+Z^2=3m^2$, and $X,Y,Z$ are odd by the
preceding modulo-eight argument. The inverse formulas become
\begin{equation}\label{eq:inverse-even}
 \begin{aligned}
 h&=(m+X+Y+Z)/2,&i&=(m+X-Y-Z)/2,\\
 j&=(m-X+Y-Z)/2,&k&=(m-X-Y+Z)/2.
 \end{aligned}
\end{equation}
They are integers, and any two coordinates differ by a sum of two odd
numbers, so the four coordinates have the same parity. Negating the triple
gives another quadruple whose corresponding coordinates sum with the first
to $m$. Since $m$ is odd, one of the two quadruples is entirely odd and the
other entirely even.

If $(X,Y,Z)$ is primitive, the odd quadruple is primitive: an odd prime
dividing all four coordinates would divide $X,Y,Z$, and two does not divide
the odd coordinates. The even quadruple is not primitive. If the normalized
triple has gcd $d>1$, then $d$ is odd, $d\mid m$ by
\cref{lem:valuation}, and \eqref{eq:inverse-even}, using that two is
invertible modulo $d$, makes both quadruples divisible by $d$.
Thus each primitive sign pair of normalized triples contributes exactly one
primitive quadruple. This proves the second line. In particular, the count
is $F(m)/2$, not a count of primitive \emph{unnormalized} triples of norm
$12m^2$.

Finally, if $4\mid n$, the number of odd coordinates of a quadruple is even
because their sum is even. Exactly two odd coordinates are impossible modulo
four, since their squares would sum to $2\pmod4$. Four odd coordinates are
impossible modulo eight, since their squares would sum to $4\pmod8$ whereas
$n^2\equiv0\pmod8$. All coordinates are therefore even, so no quadruple
is primitive.
\end{proof}

\section{An explicit lattice proof of the primitive product}\label{sec:lattice}

Write $Q(x)=x\cdot x$ on $\R^3$. A Euclidean lattice is called integral when
all its inner products are integers; its covolume is the volume of a
fundamental parallelepiped. A vector of a lattice is primitive when it is
not in $qL$ for any rational prime $q$.

\begin{lemma}[Small unimodular lattices]\label[lemma]{lem:small-lattices}
Every integral positive definite Euclidean lattice of rank $1\le r\le3$ and
covolume one has an orthonormal $\Z$-basis.
\end{lemma}
\begin{proof}
For $r=3$, the open ball of radius $4/3$ has volume
$256\pi/81>8$. Minkowski's theorem~\cite{minkowski} gives a nonzero lattice vector of
squared length less than $16/9<2$. Its squared length is a positive
integer, so it equals one. The analogous ball in rank two has area
$16\pi/9>4$, and in rank one has length $8/3>2$, so the same reasoning
applies in those ranks.

If $u\in L$ has squared length one, integrality implies
$x\cdot u\in\Z$ for every $x\in L$. Therefore
\[
 L=\Z u\mathbin{\perp}L',\qquad L'=L\cap u^\perp,
\]
because $x-(x\cdot u)u\in L'$. Orthogonal splitting preserves covolume
when the first factor has length one; thus $L'$ is integral of covolume
one. Induction on rank completes the proof.
\end{proof}

\begin{lemma}[Isotropic lines over an odd prime field]\label[lemma]{lem:isotropic-lines}
Let $p$ be an odd prime and let $\mathcal C$ be the set of isotropic
one-dimensional subspaces of $(\F_p^3,Q)$. Then $|\mathcal C|=p+1$.
For nonzero $\bar x\in\F_p^3$, the number of $\ell\in\mathcal C$ with
$\ell\subseteq\bar x^\perp$ is
\[
 1+\left(\frac{-Q(\bar x)}p\right).
\]
For $\bar x=0$, that number is $p+1$.
\end{lemma}
\begin{proof}
There is a nonzero isotropic vector: the sets of squares and of numbers
$-1-y^2$ each have $(p+1)/2$ elements and therefore intersect. This gives
a vector $(x,y,1)$ of norm zero.

Fix a nonzero isotropic vector $u$. Nondegeneracy gives $v$ with
$u\cdot v=1$, and $w=v-Q(v)u/2$ then satisfies
$Q(w)=0$ and $u\cdot w=1$. The span of $u,w$ is nondegenerate;
its one-dimensional orthogonal complement has a generator $e$ with
$Q(e)=c\ne0$. Thus the form in this basis is $2ab+ct^2$.
Replacing $a$ by $a'=-2a/c$ turns its zero equation into $a'b=t^2$.
This conic is parametrized bijectively by
$[u:v]\mapsto[u^2:v^2:uv]$. Explicitly, when its first coordinate
is nonzero, rescale to $(1,s^2,s)$; its only point with first
coordinate zero is $[0:1:0]$. It therefore has $p+1$ points.

If $Q(\bar x)=0$, use the same hyperbolic basis with $u=\bar x$.
Its orthogonal plane is spanned by $u,e$, with norm $ct^2$ on
$au+te$. Thus its sole isotropic line is $\F_p\bar x$.

If $Q(\bar x)\ne0$, the orthogonal plane is nondegenerate. Its Gram
determinant has square class $Q(\bar x)^{-1}$: adjoining $\bar x$ to
any basis of the plane gives Gram determinant $Q(\bar x)$ times the
plane determinant, whereas change of basis in the standard ternary
space changes determinant by a square. A nondegenerate binary form
has two isotropic lines when the negative of its determinant is a
square, and no isotropic lines otherwise. To see this explicitly,
diagonalize it as $As^2+Bt^2$, with $A,B\ne0$ (characteristic is
odd). An isotropic line has $t\ne0$, and $s/t$ must be a root of
$X^2=-B/A$. There are two or zero roots according as the negative
determinant $-AB$ is a square or a nonsquare. The square classes of
$-Q(\bar x)^{-1}$ and $-Q(\bar x)$ coincide. The asserted expression
follows in both cases.
\end{proof}

\begin{lemma}[The neighbors of $\Z^3$]\label[lemma]{lem:neighbors}
Put $L=\Z^3$. For each $\ell\in\mathcal C$ there is an integral
covolume-one lattice $L_\ell$ satisfying
\[
 \begin{gathered}
 L\cap L_\ell=M_\ell
    :=\{x\in L:x\bmod p\in\ell^\perp\},\\
 [L:M_\ell]=[L_\ell:M_\ell]=p,\\
 pL\subseteq M_\ell,\qquad pL_\ell\subseteq M_\ell\subseteq L.
 \end{gathered}
\]
If $v\in L\setminus pL$ and $p^2\mid Q(v)$, then $v/p$ belongs
to exactly one of these lattices, namely $L_{\F_p(v\bmod p)}$.
\end{lemma}
\begin{proof}
Choose a lift $u\in L$ of a nonzero vector of $\ell$. It can be chosen
with $p^2\mid Q(u)$. Indeed, starting with any lift, replacing it by
$u+pz$ changes $Q(u)/p$ modulo $p$ by $2u\cdot z$. The linear
functional $z\mapsto u\cdot z\pmod p$ is nonzero, hence surjective,
so a correction $z$ with the required value exists.

Set
\[
 L_\ell=M_\ell+\Z\frac up.
\]
All inner products are integral, since $u\cdot M_\ell\subseteq p\Z$
and $Q(u)\in p^2\Z$. The functional defining $M_\ell$ is surjective
modulo $p$, so its index in $L$ is $p$. The class of $u/p$ modulo
$M_\ell$ has order $p$, since $u\notin pL$ and $u\in M_\ell$.
Consequently $L_\ell$ has covolume one. If $m+tu/p\in L$, then
$p\mid t$, proving the intersection assertion. All other inclusions
follow from the definition. Each $L_\ell$ is isometric as a lattice
to $\Z^3$ by the preceding small-lattice lemma.

For the last assertion, membership forces the line of $v\bmod p$
to equal $\ell$. Conversely write $v=tu+pz$ with $p\nmid t$.
The divisibility of $Q(v)$ and $Q(u)$ by $p^2$ gives
$2t(u\cdot z)\equiv0\pmod p$. Thus $z\in M_\ell$ and
$v/p=tu/p+z\in L_\ell$.
\end{proof}

The construction is independent of the normalized lift: if another lift
is $tu+pz$ with $t\not\equiv0\pmod p$ and norm divisible by $p^2$,
the calculation in the last paragraph shows $z\in M_\ell$; the two
generators then generate the same cyclic extension of $M_\ell$.

\paragraph{Primitivity away from $p$.}
In the double count below, passage between $L$ and $L_\ell$ preserves
primitivity at every prime $q\ne p$. This can be checked without using
localization. A common vector $x\in M_\ell$ is in $qL$ if and only
if it is in $qL_\ell$. Indeed, if $x=qz$ with $z\in L$, reduction
modulo $p$ shows $z\in M_\ell$. If instead $z\in L_\ell$, its
class in the group $L_\ell/M_\ell$ of order $p$ is annihilated by
$q$ and hence is zero; again $z\in M_\ell$. For an
outside vector $y=v/p\in L_\ell\setminus L$, the corresponding
statement is $y\in qL_\ell$ if and only if $v\in qL$. In the latter
direction, write $v=qz$. Since $v\in M_\ell$ and $q$ is invertible
modulo $p$, $z\in M_\ell$. Choose integers $a,b$ with $aq+bp=1$;
then $z/p=ay+bz\in L_\ell$. The converse follows by multiplying by
$p$, since $pL_\ell\subseteq L$.

Recall that $F(m)$ is the number of primitive vectors of norm $3m^2$ in
$L$. These sets are finite by coordinate bounds. All the lattices
$L_\ell$ have exactly $F(m)$ such vectors, because they are
isometric to $L$ by a bijective $\Z$-linear isometry. Let
$\chi_p=\bigl(\frac{-3}{p}\bigr)$.

\begin{lemma}[Primitive double count]\label[lemma]{lem:double-count}
Let $p$ be an odd prime and let $m$ be positive and odd.
If $p\nmid m$, then
\[
 F(pm)=(p-\chi_p)F(m).
\]
If $p\mid m$, then
\[
 F(pm)=(p+1)F(m)-
 \left(p-\left(\frac{-3(m/p)^2}{p}\right)\right)F(m/p).
\]
\end{lemma}
\begin{proof}
Count pairs $(\ell,y)$ where $y$ is primitive in $L_\ell$ and has
norm $3m^2$. Their total number is $(p+1)F(m)$.

First count the pairs with $y\notin L$. The vector $v=py$ is in
$L\setminus pL$ and has norm $3p^2m^2$. Conversely every primitive
such $v$ gives exactly one outside pair by the neighbor lemma.
Primitivity away from $p$ is preserved as just proved. An outside
$y$ is automatically primitive at $p$, because $y\in pL_\ell$
would imply $v\in p^2L_\ell\subseteq pL$. Thus the outside pairs
number exactly $F(pm)$.

Suppose first that $p\nmid m$. The integer $3m^2$ is not divisible
by $p^2$, including at $p=3$. Therefore every vector of this norm
in an integral lattice is primitive at $p$. Inside pairs are precisely
primitive $x\in L$ with $\ell\subseteq(x\bmod p)^\perp$.
The finite-field lemma gives $1+\chi_p$ choices of $\ell$ for each
such $x$: for $p\ne3$ the factor $m^2$ is a nonzero square, and for
$p=3$ both Legendre symbols are zero. Therefore
$(p+1)F(m)=F(pm)+(1+\chi_p)F(m)$, giving the first formula.

Now suppose $p\mid m$. An inside vector that is primitive away
from $p$ has the unique form $x=p^az$ with $a\ge0$ and $z$
primitive in $L$. If $a=0$, then $x\bmod p$ is nonzero and
isotropic. Its unique orthogonal isotropic line is its own line;
for this neighbor $x/p\in L_\ell$ by the last part of the neighbor
lemma. Hence $x$ is not primitive in that neighbor and contributes
nothing. If $a\ge2$, then $x/p\in pL\subseteq L_\ell$ for every
neighbor, so again it contributes nothing.

Only $a=1$ contributes. Write $x=pz$, with $z$ primitive of norm
$3(m/p)^2$. The vector $x$ belongs to all $p+1$ neighbors. It fails
to be primitive at $p$ precisely when $z\in L_\ell$, or equivalently
when $\ell\subseteq(z\bmod p)^\perp$. The finite-field lemma shows
that the number of remaining choices is
$p-\bigl(\frac{-3(m/p)^2}{p}\bigr)$. Counting these pairs yields
the second formula.
\end{proof}

\begin{theorem}[Primitive three-square product]\label[theorem]{thm:three-squares}
For every positive odd $m$,
\[
 F(m)=8m\prod_{p\mid m}\left(1-\frac{\chi_p}{p}\right).
\]
\end{theorem}
\begin{proof}
First $F(1)=8$, since the vectors of norm three are precisely
$(\pm1,\pm1,\pm1)$.

Fix an odd prime $p$ and a positive odd $b$ with $p\nmid b$.
The first double-counting identity gives
$F(pb)=(p-\chi_p)F(b)$. Apply the second identity at $m=pb$.
Its Legendre symbol is $\chi_p$, so
\[
 F(p^2b)=(p+1)(p-\chi_p)F(b)-(p-\chi_p)F(b)
        =p(p-\chi_p)F(b).
\]
For exponent $a\ge2$, the same identity at $m=p^ab$ has zero
Legendre symbol, since $p\mid p^{a-1}b$. Induction therefore gives
\[
 F(p^{a+1}b)=(p+1)F(p^ab)-pF(p^{a-1}b)
            =p^a(p-\chi_p)F(b).
\]
Removing the prime-power factors of $m$ one at a time now gives
\[
 F(m)=8\prod_{p^a\parallel m}p^{a-1}(p-\chi_p)
     =8m\prod_{p\mid m}(1-\chi_p/p).
\]
The local symbol agrees with the mod-six character by quadratic
reciprocity, as proved explicitly in Section~\ref{sec:character}.
\end{proof}
\paragraph{The prime three.}
No use has been made of $p\nmid3$ in the neighbor argument. The
underlying Gram matrix of $L$ is the identity and stays nondegenerate
modulo every odd prime, including three. If $3\nmid m$, the norm
$3m^2$ has valuation exactly one, so inside vectors are still
primitive at three and the isotropic-line count is one. This gives
$F(3m)=3F(m)$. If $3\mid m$, the symbol in the second recurrence is
always zero; it gives
$F(3m)=4F(m)-3F(m/3)$, whose induction again yields $3F(m)$.
Thus the prime three contributes its full power through $m$, while
its Euler factor is one.



\section{Identification of the local character}\label{sec:character}

\begin{lemma}\label[lemma]{lem:character}
For every odd prime $p$,
\[
 \chi_p=\leg{-3}{p}=
 \begin{cases}
 0,&p=3,\\
 1,&p\equiv1\pmod6,\\
 -1,&p\equiv5\pmod6.
 \end{cases}
\]
\end{lemma}
\begin{proof}
The value at three is zero by definition. If $p\neq3$ is odd, quadratic
reciprocity and its supplementary law~\cite{gauss} give explicitly
\[
 \leg{-1}{p}=(-1)^{(p-1)/2},\qquad
 \leg{3}{p}\leg{p}{3}=(-1)^{(p-1)/2}.
\]
Since $\leg{p}{3}\in\{1,-1\}$, multiplication yields
$\leg{-3}{p}=\leg{p}{3}$. The nonzero squares modulo three consist only of
one. Thus this symbol is one when $p\equiv1\pmod3$ and minus one when
$p\equiv2\pmod3$. An odd prime different from three is $1$ or $5$ modulo
six, giving the displayed cases.
\end{proof}

\section{Divisibility, multiplicativity, and prime-power formulas}

\begin{theorem}[The primitive arithmetic function]\label[theorem]{thm:primitive}
Write $n=2^\alpha m$, with $m$ odd. Then $\pp(n)$ is divisible by four,
and the integer $P(n)=\pp(n)/4$ is given by
\begin{equation}\label{eq:P-product}
 P(n)=
 \begin{cases}
 \displaystyle\prod_{p^a\parallel m}p^{a-1}(p-\chi_p),&\alpha=0\text{ or }1,\\[2mm]
 0,&\alpha\geq2.
 \end{cases}
\end{equation}
In particular $P(1)=1$, $P$ is multiplicative, and, for $a\geq1$,
\begin{equation}\label{eq:P-local}
 \begin{aligned}
 P(2)&=1,&P(2^a)&=0&& (a\geq2),\\
 P(3^a)&=3^a,\\
 P(p^a)&=p^{a-1}(p-1)&&&(p\equiv1\pmod6),\\
 P(p^a)&=p^{a-1}(p+1)&&&(p\equiv5\pmod6).
 \end{aligned}
\end{equation}
\end{theorem}
\begin{proof}
The product theorem in \cref{thm:three-squares} can be written
\[
 F(m)=8\prod_{p^a\parallel m}p^{a-1}(p-\chi_p).
\]
Its right side is eight times an integer, including for $m=1$, where the
product is empty. The bridge \eqref{eq:bridge} therefore proves divisibility
of $\pp(n)$ by four and gives \eqref{eq:P-product}. Coprime positive
integers have disjoint prime factors. Their odd parts therefore give a
product of disjoint local factors, and their powers of two multiply according
to the local function with values $1,1,0,0,\ldots$. At most one of two
coprime integers is even, so this local function is multiplicative as well.
This proves multiplicativity of $P$, with $P(1)=1$. The local formulas follow
from \cref{lem:character}, including $p=3$ where $\chi_3=0$.
\end{proof}

\begin{lemma}[Content decomposition]\label[lemma]{lem:content}
For every $n\geq1$,
\begin{equation}\label{eq:raw-divisor}
 \ttq(n)=\sum_{d\mid n}\pp(d).
\end{equation}
\end{lemma}
\begin{proof}
A quadruple for positive $n$ is not the zero quadruple, so its content $c$ is
positive. Since $c$ divides the sum of its coordinates, it divides $n$.
Dividing all four coordinates by $c$ gives a primitive quadruple for $n/c$.
Conversely, scaling a primitive quadruple for a divisor $d\mid n$ by $n/d$
gives a quadruple of content exactly $n/d$. These constructions are mutually
inverse, and distinct contents give disjoint sets. Counting proves the formula.
\end{proof}

\begin{theorem}[The total arithmetic function]\label[theorem]{thm:total}
For every $n\geq1$, the raw count $\ttq(n)$ is divisible by four, and
\begin{equation}\label{eq:T-divisor}
 T(n)=\sum_{d\mid n}P(d).
\end{equation}
The function $T$ is multiplicative and, for every prime $p$ and $a\geq1$,
\begin{equation}\label{eq:T-local}
 \begin{aligned}
 T(2^a)&=2,\\
 T(3^a)&=(3^{a+1}-1)/2,\\
 T(p^a)&=p^a&& (p\equiv1\pmod6),\\
 T(p^a)&=p^a+2\frac{p^a-1}{p-1}&& (p\equiv5\pmod6).
 \end{aligned}
\end{equation}
\end{theorem}
\begin{proof}
Every summand in \eqref{eq:raw-divisor} is divisible by four by
\cref{thm:primitive}. Hence the raw total count is divisible by four, and
division of this finite sum proves \eqref{eq:T-divisor} as an identity of
integers. There is no interchange of a rounded division with a sum.

If $a,b$ are coprime, every divisor of $ab$ factors uniquely as $d_1d_2$,
with $d_1\mid a$ and $d_2\mid b$. Multiplicativity of $P$ gives
\[
 T(ab)=\sum_{d_1\mid a}\sum_{d_2\mid b}P(d_1d_2)
       =\Bigl(\sum_{d_1\mid a}P(d_1)\Bigr)
        \Bigl(\sum_{d_2\mid b}P(d_2)\Bigr)=T(a)T(b).
\]
Also $T(1)=P(1)=1$. For a prime power,
$T(p^a)=\sum_{j=0}^aP(p^j)$. At two the only nonzero summands are
$P(1)=P(2)=1$. At three the sum is $1+3+\cdots+3^a$. If
$p\equiv1\pmod6$, the sum is
$1+(p-1)\sum_{j=0}^{a-1}p^j=p^a$; if $p\equiv5\pmod6$, it is
$1+(p+1)\sum_{j=0}^{a-1}p^j=p^a+2(p^a-1)/(p-1)$.
These finite geometric sums give \eqref{eq:T-local}.
\end{proof}

\section{The equality cases}

\begin{theorem}[Exact classification]\label[theorem]{thm:equality}
For every positive integer $n$,
\begin{equation}\label{eq:equality}
 T(n)=n
 \quad\Longleftrightarrow\quad
 n=2^\varepsilon m,\quad \varepsilon\in\{0,1\},\quad m\text{ odd},
 \quad p\mid m\Longrightarrow p\equiv1\pmod6
\end{equation}
where the last condition ranges over prime divisors. The empty factorization
of $m=1$ is allowed, so both $n=1$ and $n=2$ are equality cases.
\end{theorem}
\begin{proof}
Write $n=2^\alpha m$ with $m$ odd. First, for every odd prime $p$ and every
$a\geq1$, the local formulas imply $T(p^a)\equiv1\pmod3$. For $p=3$ this
follows from the geometric sum. For $p\equiv1\pmod6$ it follows from
$p^a\equiv1\pmod3$. For $p\equiv5\pmod6$, write
\[
 T(p^a)=1+(p+1)(1+p+\cdots+p^{a-1});
\]
here $p+1$ is divisible by three. Thus $T(m)\equiv1\pmod3$ by
multiplicativity. Since $T(2^\alpha)$ is one for $\alpha=0$ and two for
$\alpha\geq1$, $T(n)$ is never divisible by three. Consequently equality
is impossible if $3\mid m$.

Assume $3\nmid m$. Every local factor $T(p^a)$ for $p\mid m$ is odd:
it is either $p^a$ or $p^a$ plus an even integer. Hence $T(m)$ is odd.
If $\alpha\geq1$, then $T(n)=2T(m)$ has exact two-adic valuation one,
whereas $n$ has valuation $\alpha$. Equality therefore forces $\alpha=1$.
This rules out compensating a higher power of two by larger odd-prime factors.

For $\alpha=0$ or $1$, equality is equivalent to $T(m)=m$. By
\eqref{eq:T-local}, every local ratio $T(p^a)/p^a$ for $p\mid m$ is at
least one, with equality exactly when $p\equiv1\pmod6$; the ratio is
strictly larger for $p\equiv5\pmod6$. The product of these positive ratios
equals one precisely when every ratio does. This proves necessity and
sufficiency in \eqref{eq:equality}, including $m=1$.
\end{proof}

\section{A finite sieve proving density zero}\label{sec:density}

Let $\B$ be the primes congruent to five modulo six. Call a positive integer
$q$ \emph{admissible} when every prime divisor of $q$ belongs to $\B$.
Such a $q$ is coprime to six. Define
\[
 a(q)=\frac{\varphi(q)}q,\qquad b(q)=\frac q{\varphi(q)}.
\]
Both are positive and $a(q)b(q)=1$, since $q\geq1$ gives
$\varphi(q)\geq1$.

\begin{lemma}[Finite descent of the sieve density]\label[lemma]{lem:sieve-step}
For every admissible $q$, there is an explicitly finite admissible extension
$q'$ with $b(q')\geq b(q)+1/6$.
\end{lemma}
\begin{proof}
Define finite sets
\[
 \begin{aligned}
 C(q)&=\{n:1\leq n\leq6q,\ n\equiv5\pmod6,\ \gcd(n,q)=1\},\\
 G(q)&=\{p\in\B:p\leq6q,\ p\nmid q\}.
 \end{aligned}
\]
There are exactly $\varphi(q)$ elements of $C(q)$. Indeed, for each
invertible residue $r$ modulo $q$, the Chinese remainder theorem, using
$\gcd(6,q)=1$, gives exactly one residue modulo $6q$ with residue five
modulo six and residue $r$ modulo $q$. Its representative in $[1,6q]$ is
positive and coprime to $q$; this parametrizes $C(q)$ bijectively.

Every positive integer $n\equiv5\pmod6$ has a prime divisor in $\B$.
It is divisible by neither two nor three, so all its prime factors are one
or five modulo six. If all were one, their product with multiplicities would
be one modulo six, a contradiction. For $n\in C(q)$, such a divisor cannot
divide $q$ and is at most $n\leq6q$. It therefore belongs to $G(q)$.

The union bound and the exact count of positive multiples of $p$ give
\[
 \varphi(q)=|C(q)|
 \leq\sum_{p\in G(q)}\left\lfloor\frac{6q}p\right\rfloor
 \leq6q\sum_{p\in G(q)}\frac1p.
\]
Consequently
\begin{equation}\label{eq:finite-reciprocal-bound}
 \sum_{p\in G(q)}\frac1p\geq\frac{a(q)}6.
\end{equation}
Set $q'=q\prod_{p\in G(q)}p$. The new primes are distinct and coprime
to $q$, so $q'$ is admissible and multiplicativity of the totient gives
\[
 b(q')=b(q)\prod_{p\in G(q)}\frac p{p-1}.
\]
Each factor is at least $1+1/p$. For nonnegative numbers $x_p$, expanding
a finite product gives $\prod(1+x_p)\geq1+\sum x_p$. Using
\eqref{eq:finite-reciprocal-bound}, we conclude
\[
 b(q')\geq b(q)\left(1+\sum_{p\in G(q)}\frac1p\right)
 \geq b(q)\left(1+\frac{a(q)}6\right)=b(q)+\frac16.
\]
All prime sets in this construction are bounded by $6q$ and finite.
\end{proof}

\begin{corollary}[Arbitrarily small finite sieve products]\label[corollary]{cor:small-sieve}
For every real $\delta>0$, an admissible $q$ satisfies $\varphi(q)/q<\delta$.
\end{corollary}
\begin{proof}
Starting with $q_0=1$, apply \cref{lem:sieve-step} repeatedly. Each step
constructs a finite positive integer, and induction gives
\[
 b(q_t)\geq1+t/6,\qquad a(q_t)\leq\frac6{t+6}.
\]
Choose an integer $t>6/\delta$. Then $a(q_t)<\delta$.
\end{proof}

\begin{theorem}[Elementary density-zero sieve]\label[theorem]{thm:sieve-density}
Let $E$ be a set of positive integers, none of which has a prime divisor
congruent to five modulo six. Then
\[
 \lim_{N\to\infty}\frac{|E\cap[1,N]|}N=0.
\]
\end{theorem}
\begin{proof}
Every member of $E$ is coprime to every admissible $q$. Each complete block
of $q$ consecutive integers contains exactly $\varphi(q)$ integers coprime
to $q$. The final incomplete block contains at most that many. Therefore
\[
 |E\cap[1,N]|\leq\varphi(q)\bigl(\lfloor N/q\rfloor+1\bigr)
 \leq\frac{\varphi(q)}qN+\varphi(q).
\]
Given $\epsilon>0$, use \cref{cor:small-sieve} to choose an admissible $q$
with $\varphi(q)/q<\epsilon/2$. For every integer
$N>2\varphi(q)/\epsilon$, division of the preceding inequality by $N$
gives
\[
 0\leq\frac{|E\cap[1,N]|}N<\epsilon.
\]
This proves the asserted limit; existence of the density was not assumed.
\end{proof}

\begin{corollary}[Density of the equality set]\label[corollary]{cor:equality-density}
The set $\{n\geq1:T(n)=n\}$ has natural density zero.
\end{corollary}
\begin{proof}
By \cref{thm:equality}, every such $n$ has no prime divisor congruent to
five modulo six. Apply \cref{thm:sieve-density}.
\end{proof}

In particular, the positive-asymptotic-density conjecture recorded in
A354766 is false: the density exists and equals zero.

\section{Verification and formalization boundary}

The representation product in \cref{thm:three-squares} is proved here from
the explicitly constructed neighbors, their elementary low-rank lattice
classification, and a double count. The published
Cooper--Hirschhorn formula agrees with it, including the prime-three factor,
but is only a cross-check. The density argument is a finite construction
followed by an explicit epsilon bound; reciprocal-prime divergence is not
one of its assumptions.

The recovered verification program independently enumerates the original
quadruples and the Hadamard inverse triples for $1\leq n\leq150$. It checks
the corrected bridge, primitive and total counts, local product formulas,
and divisibility of both raw counts by four. Those checks passed after
recovery, and the resulting CSV agrees with the saved historical audit.
An independent cloud computation additionally checked both defining
quadruple models through $n=300$ and primitive three-square counts at all
38 odd levels through $75$, with no discrepancy.
Finite computations are consistency checks, not proof steps in this paper.

The companion Lean 4.27.0 development implements the lattice
classification, normalized neighbors, finite conic counts, primitive-vector
double count, prime-power recurrences, and Euler product. It also verifies
the primitive A278085 complement bijection and the finite density sieve.
The default dependency graph has 43 local modules with no proof admissions
or added axioms. All 14 concluding dependency reports contain only
\texttt{propext}, \texttt{Classical.choice}, and \texttt{Quot.sound}; neither
an assumed representation formula nor native-compiler reduction is a
dependency of those conclusions. The sources pin Mathlib commit
\texttt{a3a10db0e9d66acbebf76c5e6a135066525ac900} and include the checks.

The mathematical derivations have undergone internal adversarial checking.
This document remains a draft for independent human mathematical review.
AI systems assisted exploration, checking, and drafting; that assistance
is disclosed and is not itself evidence of correctness.

\begin{thebibliography}{9}

\bibitem{oeis-total}
OEIS Foundation Inc., \emph{A354766},
\url{https://oeis.org/A354766}.

\bibitem{oeis-primitive}
OEIS Foundation Inc., \emph{A278085},
\url{https://oeis.org/A278085}.

\bibitem{minkowski}
H. Minkowski, \emph{Geometrie der Zahlen}, B. G. Teubner, Leipzig, 1896.
The centrally symmetric convex-body lattice-point theorem.

\bibitem{gauss}
C. F. Gauss, \emph{Disquisitiones Arithmeticae}, Leipzig, 1801,
Article 131 (quadratic reciprocity), with the supplementary law for $-1$.

\bibitem{kneser}
M. Kneser, \emph{Klassenzahlen definiter quadratischer Formen},
Arch. Math. \textbf{8} (1957), 241--250.

\bibitem{voight}
J. Voight, \emph{Kneser's method of neighbors}, April 30, 2024,
Proposition 3.13 and Corollary 3.15,
\url{https://jvoight.github.io/articles/kneser-033024.pdf}.

\bibitem{cooper-hirschhorn}
S. Cooper and M. D. Hirschhorn,
\emph{On the number of primitive representations of integers as sums of squares},
Ramanujan J. \textbf{13} (2007), 7--25,
\href{https://doi.org/10.1007/s11139-006-0240-6}{doi:10.1007/s11139-006-0240-6}.

\end{thebibliography}

\end{document}
